% SKELETON (fill batch): J — inherited RTSC magnet toolchain.
% TODO (fill batch): fold the RTSC inheritance provenance — solenoid_axisym
%   .geo/.pro (GetDP 4.0) + the Wheeler on-axis closed-form anchor (M318) +
%   the 5-axis RtscVerifyRecord schema extension (device=ioffe_pritchard_trap)
%   from RTSC.md §4.2 and ~/core/hexa-lang/stdlib/rtsc/ + stdlib/material/magnet/.
\section{Inherited RTSC Magnet Toolchain}
\label{app:rtsc}

This appendix documents the cross-domain magnet inheritance that powers
process \textcircled{\scriptsize 6} (Appendix~\ref{app:confinement}). It will
trace each reused asset to its RTSC origin --- the GetDP\,4.0 axisymmetric
A--$\phi$ magnetostatic deck (\texttt{solenoid\_axisym.geo}/\texttt{.pro})
\citep{getdp}, the Wheeler on-axis closed-form current-loop anchor (the M318
cross-check), and the RTSC verify-record schema extended with
\texttt{device = ioffe\_pritchard\_trap} --- and argue why the same
magnetostatic spine serves a field \emph{well} (this antihydrogen trap) and a
field \emph{peak} (an RTSC high-field solenoid) with only a coil-geometry
change. Provenance: \texttt{RTSC.md} \S4.2, hexa-lang
\texttt{stdlib/rtsc/} and \texttt{stdlib/material/magnet/}.

% TODO (fill batch): the asset-by-asset provenance table, the peak-vs-well
% geometry contrast, the verify-once-reuse-twice SSOT argument, the deck
% derivation runbook.
